\documentclass[12pt]{article}

\usepackage{amsmath}
\usepackage{geometry}
\usepackage{hyperref}
\usepackage{xcolor}

\geometry{margin=1.2in}

\hypersetup{colorlinks=true, linkcolor=blue!60!black}

\newcommand{\TODO}[1]{\textcolor{red}{\textbf{[TODO: #1]}}}

\title{Spatial rescue in a biofilm:\\a first-principles narrative}
\author{}
\date{}

\begin{document}
\maketitle

\noindent This document explains the model, its mechanisms, and the central
question from first principles, without reference to simulation results or
formal derivations. It is meant to build intuition before reading
\texttt{derivation.tex} or \texttt{writeup.tex}.

% ---------------------------------------------------------------------------
\section{The biological question}
% ---------------------------------------------------------------------------

Bacteria in a biofilm are physically shielded from antibiotic. Drug diffuses
in from the outside, so only the outermost layer of cells is exposed to
lethal concentrations. Cells buried deeper are protected. Does that
protection actually help the population survive?

We study this through \emph{evolutionary rescue}: a declining population
avoids extinction because a resistant mutant arises and takes over before
the population collapses. The question is whether the biofilm's spatial
structure---a protected interior feeding into an antibiotic-exposed
surface---makes rescue more or less likely than in a well-mixed population
of the same size.

% ---------------------------------------------------------------------------
\section{The model}
% ---------------------------------------------------------------------------

The biofilm is a line of cells. The outermost $l$ cells are the
\textbf{edge}: exposed to antibiotic, dying faster than they reproduce. The
rest is the \textbf{core}: shielded, with births and deaths roughly balanced
(neutral).

The core and edge are not isolated. Every edge death draws the innermost
core cell outward into the edge; every edge birth pushes the innermost edge
cell inward into the core. This \textbf{conveyor belt} means the two
compartments exchange cells continuously, driven by births and deaths at the
boundary. Cells do not stay put.

Wild-type (WT) cells are killed by antibiotic in the edge. Resistant (R)
cells survive and grow there. Resistance arises by mutation at a fixed
probability per birth event. We compare this biofilm to a well-mixed
population of the same initial size, using only edge rates (every cell is
exposed), as the baseline.

% ---------------------------------------------------------------------------
\section{Two regimes of population decline}
% ---------------------------------------------------------------------------

Once antibiotic treatment starts, the edge is net-dying. The conveyor belt
means the core shrinks as it feeds replacement cells into the edge. This
produces two qualitatively different periods.

\textbf{The buffered decline} (while the core persists): Every edge death
is compensated by a core cell entering the edge, so the edge stays at
exactly $l$ cells. The total population falls, but linearly and
slowly---only the $l$ edge cells drive the net decline. The rate of decline
is proportional to dose times $l$. This period ends when the last core cell
is consumed.

\textbf{The final collapse} (after the core is exhausted): No buffer
remains. The whole remaining population is in the edge and falls
exponentially, exactly as a well-mixed population would. The biofilm has
lost its special structure.

How long the buffered decline lasts depends on dose: at low dose, the edge
declines slowly, so the core is consumed slowly and the buffered period is
long. At high dose, the core is consumed quickly.

% ---------------------------------------------------------------------------
\section{Naive intuition: the biofilm should win}
% ---------------------------------------------------------------------------

The buffered decline extends the population's lifetime. If the core is
metabolically active---cells are dividing and dying there---it also
contributes extra births. Extra births mean extra chances for resistance
to arise. The biofilm appears to have all the advantages: more time, more
mutations, a protected reservoir.

This intuition is wrong, or at least incomplete. The same spatial structure
that creates the buffer also creates two suppression mechanisms that can
hurt rescue.

% ---------------------------------------------------------------------------
\section{Suppression mechanism 1: WT resupply}
% ---------------------------------------------------------------------------

During the buffered decline, every WT death in the edge immediately draws in
a fresh WT cell from the core. As long as the core is WT-dominated (which it
is throughout most of the buffered decline), the WT count in the edge is
actively \emph{maintained} at $l$ cells. The antibiotic is killing WT at the
edge, but the core is instantly replacing every one.

A lone R mutant that arises in the edge cannot capitalize on this. In a
well-mixed population, WT cells that die are gone---the competitive
landscape for R improves as the population shrinks. In the biofilm, dying
WT cells are replaced from the core. The competitive pressure on R never
weakens. A lone R cell's birth rate equals WT's but its deaths immediately
inject a fresh WT competitor; the effective net rate for R in the edge is
negative during the buffered decline.

Two properties of this mechanism shape everything that follows.

\textbf{It is dose-independent.} The WT resupply rate scales with dose, but
so does the WT death rate that drives it---these cancel. The suppression of
R establishment is the same regardless of how hard the antibiotic is
working.

\textbf{It operates even with a dormant core.} Core cells slide into the
edge mechanically when edge cells die, regardless of whether the core cells
themselves divide. WT resupply is driven by edge deaths, not core births.

% ---------------------------------------------------------------------------
\section{Suppression mechanism 2: drift kills core mutations before delivery}
% ---------------------------------------------------------------------------

A resistant mutant born in the core is in a neutral environment---no
antibiotic there, no growth advantage. The lineage drifts stochastically
while waiting for the retreating boundary to expose it to the edge, where
it would have an advantage.

At \textbf{low dose}, the boundary retreats slowly. A mutation born deep in
the core waits a long time. Neutral drift has many opportunities to
eliminate the lineage before the boundary ever reaches it. At \textbf{high
dose}, the boundary retreats fast and sweeps up lineages before drift can
act. Core delivery is efficient at high dose and nearly useless at low dose.

This mechanism \textbf{requires an active core} ($b_c > 0$). A dormant core
produces no mutations; there is nothing to be lost to drift.

% ---------------------------------------------------------------------------
\section{The dormant core: a clean baseline}
% ---------------------------------------------------------------------------

Before combining the two mechanisms, it is worth working through the dormant
core case ($b_c = 0$) because it gives the clearest prediction.

With a dormant core, there are no core mutations. Rescue must come entirely
from the edge. Now consider how many edge birth events the biofilm edge
produces over its entire lifetime, compared to the well-mixed population.

During the buffered decline the edge is held at exactly $l$ cells,
generating births at rate $b_e \cdot l$ continuously. The buffered decline
lasts until the core is consumed---a duration proportional to $(N_0 -
l)/(s_e \cdot l)$. Multiplying: the edge produces $(N_0 - l)/s_e$ birth
events during the buffered decline. After the core is exhausted, the final
collapse contributes another $l/s_e$ edge birth events. Total:
$(N_0 - l + l)/s_e = N_0/s_e$.

The well-mixed population, declining exponentially from $N_0$, produces
$N_0/s_e$ total birth events.

\textbf{They are equal.} The biofilm's longer lifetime exactly compensates
for having only $l$ cells in the edge at any one time. A dormant core adds
no mutation supply.

But the edge mutations in the biofilm are not equal in quality to the
well-mixed mutations. Every edge mutation during the buffered decline faces
WT resupply---establishment is suppressed. The well-mixed mutations face
no such suppression. Therefore:

\begin{quote}
\textit{A dormant-core biofilm has the same total mutation supply as a
well-mixed population of the same size, but worse establishment probability
for every mutation. It rescues strictly less often, at every dose, by a
dose-independent factor.}
\end{quote}

There is no crossover for a dormant core. It is uniformly worse.

% ---------------------------------------------------------------------------
\section{The active core: a crossover with dose}
% ---------------------------------------------------------------------------

With a metabolically active core ($b_c > 0$), the picture is more complex,
and a dose-dependent crossover emerges.

\subsection*{At low dose}

Core delivery is drift-suppressed---most core-born lineages die out before
the slow-moving boundary reaches them. As dose approaches zero, core
delivery approaches zero entirely. The active-core biofilm behaves
essentially like the dormant-core biofilm at low dose: rescue comes from
edge mutations, which face the same WT resupply suppression as before.

Edge mutations in the biofilm are competing against the entire well-mixed
population. The biofilm's edge ($l$ cells) generates the same total births
as the whole well-mixed population ($N_0$ cells), but establishment is
suppressed throughout the buffered decline. Well-mixed wins.

\textbf{At low dose, the biofilm with an active core still rescues less
often than well-mixed.} The core is not yet helping---it is contributing
neither mutations (drift kills them) nor any structural advantage. The
conveyor belt is acting purely as a WT resupply machine.

\subsection*{At high dose}

Core delivery becomes efficient. The fast-retreating boundary sweeps up
core-born lineages before drift eliminates them. Moreover, a lineage that
survives long enough in the core before delivery has had time to grow: it
arrives at the edge not as a single cell but as a \emph{bolus} of several
cells. A multi-cell bolus has disproportionately better establishment
prospects than the same number of independent single-cell arrivals, because
the rescue threshold is a discrete nonlinear criterion. This bolus advantage
is not available to the well-mixed population.

At high dose, the buffered decline is also shorter, so WT resupply suppresses
edge mutations for less time before the core is exhausted.

\textbf{At high dose, the biofilm with an active core rescues more often
than well-mixed.} The core delivery pathway is now genuinely contributing,
and the bolus advantage compounds it.

\subsection*{The crossover}

A crossover exists between low and high dose---the biofilm transitions from
rescuing less often than well-mixed to more often. The crossover is driven
by whether core delivery is efficient enough to compensate for the
establishment suppression. It requires an active core; a dormant core has
no crossover.

The location of the crossover depends on:
\begin{itemize}
    \item \textbf{Core turnover rate $b_c$}: faster turnover generates more
    core mutations but also means faster neutral drift, so the two partially
    cancel. More turnover shifts the crossover to lower dose.
    \item \textbf{Edge width $l$} and \textbf{core size $N_0 - l$}: a larger
    core relative to the edge means more potential core delivery, but also a
    longer buffered decline during which WT resupply suppresses edge
    mutations. The balance determines the crossover dose.
    \item \textbf{The establishment suppression}: how much WT resupply
    actually suppresses each edge mutation's probability of establishing.
    This depends on $l$, $b_R$, $b_e$, and $d_R$, and is not yet fully
    characterized.
\end{itemize}

\subsection*{The irony}

The same antibiotic pressure that kills the population is the mechanism
that makes the core useful. High dose drives fast boundary retreat, which
is what delivers core-born mutants efficiently and what limits the duration
of WT resupply suppression. A strong antibiotic simultaneously threatens
the population and activates the biofilm's one genuine advantage.

% ---------------------------------------------------------------------------
\section{What we do not yet know}
% ---------------------------------------------------------------------------

The narrative above gives the qualitative picture from first principles, but
several things remain to be worked out before we can predict the crossover
dose or compare quantitatively to data.

\textbf{The actual establishment suppression.} The WT resupply mechanism
suppresses R establishment during the buffered decline. How severe this
suppression is depends on the spatial arrangement of R cells in the edge.
If R cells born in the edge cluster near the outer surface---which is
plausible since births occur adjacent to the parent---they may be less
vulnerable to WT-resupply displacement than a uniform-mixing model would
predict. The true suppression could be substantially weaker than the
worst-case analytical estimate.

\textbf{The crossover dose.} Combining the establishment suppression with
the core delivery efficiency, the crossover dose can in principle be
calculated. Whether it falls within a biologically relevant range, and how
it shifts with $b_c$, $l$, and other parameters, needs to be worked out.

\textbf{Whether the crossover is observable.} At doses where rescue
probability is nearly one for both biofilm and well-mixed, the probability
metric becomes uninformative. The crossover in rescue probability may be
easier to detect through time-to-rescue: even when both populations
rescue with certainty, the biofilm may take longer at low dose (edge
establishment suppressed) and shorter at high dose (bolus delivery).

\end{document}
